% §5.5 What Comes Next
%
% Monograph register: compact Part II forward paragraph. Cross-refs
% sec:from-atom-to-set, sec:bloom-from-scratch, sec:predicate-algebra,
% sec:one-sided-two-sided, ch:set-composition, ch:classification-measures,
% ch:information-bounds, ch:optimal-bernoulli-sets, ch:perfect-hash-filters,
% sec:composition-theorem, sec:sampling-rates. DEFINED: sec:ch5-what-comes-next.

\section{What Comes Next}
\label{sec:ch5-what-comes-next}

The chapter lifted the noisy bit to a function $\ASet{A} : \Universe \to
\Bernoulli{\BoolType}$ (\Cref{sec:from-atom-to-set}), built the Bloom filter
as its canonical realization (\Cref{sec:bloom-from-scratch}), placed both
inside the chapter-4 axioms so that two rates fix the entire joint behavior
(\Cref{sec:predicate-algebra}), and separated the rate shapes an approximate
set can take (\Cref{sec:one-sided-two-sided}). The rest of Part~II develops
the algebra and bounds the construction underwrites.
\Cref{ch:set-composition} derives the rates of unions, intersections,
complements, and differences as the composition theorem of
\Cref{sec:composition-theorem} applied to the predicate algebra, and asks
which subfamilies, the perfect filter line among them, are closed under
each operation. \Cref{ch:classification-measures} turns the two intrinsic
rates into the workload-dependent measures practitioners report, positive
and negative predictive value, recall, $F_1$, and supplies the proper
interval estimators that \Cref{sec:sampling-rates} deferred.
\Cref{ch:information-bounds} proves the entropy floor of $-\log_2 \fprate$
bits per element for an FNR-zero set and locates the standard Bloom filter's
$1.44\times$ overhead against it. The two constructions that take the
perfect filter line as their subject,
\Cref{ch:optimal-bernoulli-sets} closing the FNR-zero gap and
\Cref{ch:perfect-hash-filters} realizing the FPR-zero side, arrive in
Parts~IV and~V once the intervening theory is in place.
